\textbf{结论名称：}等比数列片段和性质

\textbf{适用条件：}等比数列$\{a_n\}$，公比$q\neq 1$，且所取片段和均有意义（通常要求$S_m\neq 0$）。

\textbf{变量说明：}设$\{a_n\}$首项为$a_1$，公比$q$，前$n$项和为$S_n$；取正整数$m$，记$A_1=S_m$，$A_2=S_{2m}-S_m$，$A_3=S_{3m}-S_{2m}$。

\textbf{核心公式：}
\[
\boxed{S_m,\ S_{2m}-S_m,\ S_{3m}-S_{2m},\ \ldots\ \text{成等比数列，公比为}\ q^m}
\]
即当$S_m\neq0$时，有
\[
\frac{S_{2m}-S_m}{S_m}= \frac{S_{3m}-S_{2m}}{S_{2m}-S_m}=q^m.
\]

\textbf{使用场景：}处理等比数列中“等间距分段和”的比值问题，可避免逐项求和，直接利用$q^m$求出后续片段和或判断片段和的等比关系。

\textbf{简单例子：}已知等比数列$1,2,4,8,16,32,\dots$，$a_1=1,\ q=2$。取$m=2$，则
$S_2=1+2=3$，$S_4-S_2=(1+2+4+8)-3=12$，$S_6-S_4=(1+2+4+8+16+32)-15=48$。
所得片段和$3,12,48$成等比数列，公比$4=2^2$，满足$12/3=48/12=4$。